% 原书物理页=印刷页39--41；4.3 全节.
\section{诱导表示与诱导特征标}\label{sec:fg-4-3}

我们现在来研究一种从群 $G$ 的子群 $H$ 的表示出发构造 $G$ 的表示的方法.

\begin{numpara}[特征标的观点]\label{np:fg-4-3-1}
设 $(V,\rho)$ 是 $G$ 的一个表示. 把 $\rho$ 限制到 $H$ 上，便得到 $H$ 的一个表示 $(V,\rho_{|H})$. 它的特征标同样由限制给出：
\[
\chi_{\rho|H}=(\chi_\rho)_{|H}.
\]
注意，若 $\rho$ 不可约，则 $\rho_{|H}$ 一般并不可约（试想极端情形 $H=\{e\}$）.\footnote{译注：原文此处作 ``ext\^eme''，脱一字母 r，应为 ``extr\^eme''（极端的）.}

反过来，从 $H$ 的一个表示 $(W,\tau)$ 出发. 把它的特征标的零延拓记作 $\chi_\tau^0$. 于是有
\[
\chi_\tau^0(g)=\begin{cases}\chi_\tau(g), & \text{若 } g\in H,\\ 0, & \text{若 } g\notin H.\end{cases}
\]
一般说来，这个函数在 $H$-共轭下不变，但在 $G$-共轭下则不然. 使它对 $G$-共轭不变的一种``极小''的办法是令
\[
\chi_\tau^G(g):=\sum_{x\in H\backslash G}\chi_\tau^0(xgx^{-1})=\frac{1}{|H|}\sum_{x\in G}\chi_\tau^0(xgx^{-1}).
\]
于是 $\chi_\tau^G$ 是 $G$ 上的一个中心函数，并且当 $g$ 不与 $H$ 中元素共轭时 $\chi_\tau^G(g)=0$. 在下面的定理中，把 $\mathcal{F}(G,\mathbb{C})$ 上的内积记作 $\langle.,\,.\rangle_G$，而把 $\mathcal{F}(H,\mathbb{C})$ 上的内积记作 $\langle.,\,.\rangle_H$.
\end{numpara}

\begin{theorem*}
设 $H<G$，$(W,\tau)$ 如上.
\begin{enumerate}[label=\roman*)]
\item 对 $G$ 上任意中心函数 $f$，有 $\langle\chi_\tau^G,f\rangle_G=\langle\chi_\tau,f_{|H}\rangle_H$.
\item 中心函数 $\chi_\tau^G$ 是 $G$ 的某个表示的特征标.
\end{enumerate}
\end{theorem*}
\begin{proof}
对于 i)，我们计算
\begin{align*}
\langle\chi_\tau^G,f\rangle_G
&=\frac{1}{|G||H|}\sum_{g\in G}\sum_{x\in G}\chi_\tau^0(xgx^{-1})\overline{f(g)}
=\frac{1}{|G||H|}\sum_{x\in G}\sum_{g\in G}\chi_\tau^0(xgx^{-1})\overline{f(xgx^{-1})}\\
&=\frac{1}{|G||H|}\sum_{x\in G}\sum_{g\in x^{-1}Hx}\chi_\tau(xgx^{-1})\overline{f(xgx^{-1})}
=\frac{1}{|G||H|}\sum_{x\in G}\sum_{h\in H}\chi_\tau(h)\overline{f(h)}\\
&=\frac{1}{|G|}\sum_{x\in G}\langle\chi_\tau,f_{|H}\rangle_H=\langle\chi_\tau,f_{|H}\rangle_H.
\end{align*}
特别地，对 $G$ 的任意不可约表示 $\rho$，有
\[
a_{\tau\rho}:=\langle\chi_\tau^G,\chi_\rho\rangle_G=\langle\chi_\tau,\chi_{\rho|H}\rangle_H\in\mathbb{N}.
\]
由等式 $\chi_\tau^G=\sum_\rho a_{\tau\rho}\chi_\rho$ 可知，$\chi_\tau^G$ 是表示 $\bigoplus_\rho\rho^{\oplus a_{\tau\rho}}$ 的特征标.
\end{proof}

特征标 $\chi_\tau^G$ 称为 $\chi_\tau$ 到 $G$ 的\newterm{fg-induced-character}{诱导特征标}{induced character}.

\noindent\textbf{例.}
\begin{enumerate}[label=\roman*)]
\item 正则表示 $\mathbb{C}G$ 的特征标是平凡子群 $H=\{e\}$ 的平凡表示 $\mathbf{1}$ 的诱导特征标. 这可以写作 $\chi_{\mathbb{C}G}=\chi^G_{\mathbf{1}_{\{e\}}}$.
\item 更一般地，设 $X$ 是形如 $G/H$ 的传递 $G$-集，则置换表示 $\mathbb{C}X$ 的特征标是 $H$ 的平凡表示 $\mathbf{1}_H$ 的诱导特征标，即 $\chi_{\mathbb{C}X}=\chi^G_{\mathbf{1}_H}$.
\item 设 $G=D_n=\langle r\rangle\rtimes\langle s\rangle$. 前面已经看到，$G$ 的标准特征标 $\chi_{\mathrm{std}}$ 在循环子群 $\langle r\rangle$ 之外恒为零，且由下式给出：
\[
\chi_{\mathrm{std}}(r^k)=2\cos(2\pi k/n)=e^{2i\pi k/n}+e^{-2i\pi k/n}
\]
对一切 $k$ 成立. 设 $\psi$ 是群 $\langle r\rangle$ 的（维数为 $1$ 的不可约）特征标，由 $\psi(r^k):=e^{2i\pi k/n}$ 给出. 由于 $sr^ks^{-1}=r^{-k}$，可见 $\chi_{\mathrm{std}}$ 就是 $\psi$ 的诱导特征标，即 $\chi_{\mathrm{std}}=\psi^G$.
\end{enumerate}

第三个例子促使我们去寻找一个判别准则，使得一个不可约特征标的诱导特征标仍然是不可约的. 那样我们便拥有了一个制造不可约特征标的实用工具.

为简单起见，设 $H$ 是 $G$ 的正规子群，并设 $\chi=\chi_\tau$ 是 $H$ 的一个不可约特征标. 对每个 $x\in G$，函数 $\chi^x:h\mapsto\chi(xhx^{-1})$ 是 $H$ 上的中心函数. 它是表示 $(W,\tau^x)$ 的特征标，其中 $\tau^x(h):=\tau(xhx^{-1})$；与 $(W,\tau)$ 一样，这个表示也是不可约的. 特征标 $\chi^x$ 称为 $\chi$ 被 $x$ 作用所得的\newterm{fg-conjugate-character}{共轭特征标}{conjugate character}.

\begin{proposition*}
假设对一切 $x\in G\setminus H$ 都有 $\chi\neq\chi^x$.\footnote{译注：原句首词印作 ``Suppposons''，衍一字母 p，应为 ``Supposons''（假设）.} 那么诱导特征标 $\chi^G$ 是不可约的.
\end{proposition*}
\begin{proof}
只需验证 $\langle\chi^G,\chi^G\rangle_G=1$. 根据前面的定理，有 $\langle\chi^G,\chi^G\rangle_G=\langle\chi,(\chi^G)_{|H}\rangle_H$. 而由定义，$(\chi^G)_{|H}=\sum_{x\in G/H}\chi^x$. 于是 $\langle\chi^G,\chi^G\rangle_G=\sum_{x\in G/H}\langle\chi,\chi^x\rangle$，由我们的假设与特征标的正交性，它等于 $1$.
\end{proof}

\begin{numpara}[表示的观点]\label{np:fg-4-3-2}
与前面一样，从 $G$ 的子群 $H$ 的一个表示 $(W,\tau)$ 出发. 令
\[
\operatorname{Ind}_H^G(W):=\{f:G\longrightarrow W,\ \forall h\in H,\ \forall x\in G,\ f(hx)=\tau(h)f(x)\}.
\]
我们在这个向量空间上赋予右平移作用 $(gf)(x):=f(xg)$. 这样便得到一个表示，记作 $\operatorname{Ind}_H^G(\tau)$.
\end{numpara}

\begin{exercise*}
证明 $\operatorname{Ind}_H^G(\tau)$ 的特征标等于诱导特征标 $\chi_\tau^G$.
\end{exercise*}

表示的观点比起特征标的观点，其优点在于它对任意系数域 $K$ 都适用，甚至是正特征的.
